% TOP_PROBLEM: {"id": 5300057, "title": "Boundary theorems for geometric coding trees", "category_id": 11, "category": {"id": 11, "name": "dynamical_systems", "display_name": "Dynamical Systems", "description": "Problems about long-term behavior of deterministic systems, Hamiltonian dynamics, and periodic orbits.", "slug": "dynamical-systems", "order_index": 11, "created_at": "2026-07-31T15:26:25.671Z"}, "status": "open"}
% FIRST_POSTED: null
% ENTRY_KIND: statement_only
% Imported from: batch05.tex; UnsolvedMath ID 5300057
\documentclass[11pt]{article}
\usepackage[a4paper,margin=25mm]{geometry}
\usepackage{amsmath,amssymb,mathtools}
\usepackage{fontspec,unicode-math}
\setmainfont{Latin Modern Roman}
\setsansfont{Latin Modern Sans}
\setmathfont{Latin Modern Math}
\usepackage{xurl,hyperref}
\hypersetup{colorlinks=true,linkcolor=blue,urlcolor=blue,pdftitle={UnsolvedMath Problem 5300057}}
\newcommand{\sref}[2]{\hyperlink{source:#1:#2}{[S#2]}}
\newcommand{\cataloguescope}[1]{\paragraph{Mathematical question.}#1}
\begin{document}
% UnsolvedMath ID 5300057; source catalogue code AMR-052-0057
\setcounter{section}{5300056}
\section{Boundary theorems for geometric coding trees}\label{5300057}
{\small\sffamily UnsolvedMath ID 5300057 · Dynamical Systems}
\subsection{Definitions and mathematical statement}

\paragraph{Shared notation.}$\mathbb N=\{1,2,\ldots\}$, and $\mathbb Z,\mathbb Q,\mathbb R,\mathbb C$ denote the integers, rationals, reals and complex numbers. Any other conventions are specified locally.


Let $U\subset\widehat{\mathbb C}=\mathbb C\cup\{\infty\}$ be open and $f:U\to f(U)$ holomorphic and proper (inverse images of compact sets are compact), with $f(U)\supset U$. Choose a root $z\in f(U)$, $d\geq2$ distinct preimages $z^1,\ldots,z^d\in U$, and curves $\gamma^i\subset f(U)$ from $z$ to $z^i$. Here $\mathbb N_0=\{0,1,2,\ldots\}$ and $\Sigma_d=\{1,\ldots,d\}^{\mathbb N_0}$, with left shift $\sigma$. A geometric coding tree consists of choices of successive lifted curves $\gamma_n(\alpha)$ satisfying $\gamma_0(\alpha)=\gamma^{\alpha_0}$, $f\circ\gamma_{n+1}(\alpha)=\gamma_n(\sigma\alpha)$ and $\gamma_{n+1}(\alpha)(0)=\gamma_n(\alpha)(1)$. Write $z_n(\alpha)=\gamma_n(\alpha)(1)$, $D(z_\infty)=\{\alpha:z_n(\alpha)\text{ converges in }\overline U\}$ and $z_\infty(\alpha)=\lim_n z_n(\alpha)$. The construction does not impose absence of critical values when suitable lifts can be chosen. A Riemann map is a biholomorphic map $R:\mathbb D\to V$ onto a simply connected proper plane domain, where $\mathbb D=\{w:|w|<1\}$. Its radial boundary values are limits of $R(r\zeta)$ as $r\uparrow1$ for $|\zeta|=1$. In the coding-tree dictionary, a prime end corresponds to a geometric branch, its impression to the intersection of the closures of images of longer and longer symbolic cylinders, and its principal points to limit points of the branch vertices.
\paragraph{Statement recorded in the supplied source.}
Which theorems about boundary behaviour of Riemann maps have analogues for geometric coding trees and their coding maps $z_\infty$? \sref{5300057}{2}
\subsection{Short English statement}
Determine which results about approaching the boundary through a Riemann map also hold when the approach is through branches of a geometric coding tree.
\subsection{Sources}
\begin{description}
\item[\hypertarget{source:5300057:1}{S1}] ULAM.AI and the UnsolvedMath Contributors, UnsolvedMath: A Curated Collection of Open Mathematics Problems, record 5300057 (AMR-052-0057). CC BY 4.0. \url{https://huggingface.co/datasets/ulamai/UnsolvedMath}.
\item[\hypertarget{source:5300057:2}{S2}] Feliks Przytycki, ``On Invariant Measures for Iterations of Holomorphic Maps'', in Ben Bielefeld and Mikhail Lyubich (eds.), Problems in Holomorphic Dynamics (1992), Problem 2.6, pp. 33; arXiv:math/9205209v1. \url{https://arxiv.org/pdf/math/9205209}.
\end{description}
\label{end:5300057}
\end{document}
